\section{Data Model}
\label{sec:datamodel}

The schema is small (ten tables) and intentionally normalised only as far as
the access patterns require. Figure~\ref{fig:db-erd} shows the full ERD;
this section walks each table and the rationale.

\fig{db-erd}{0.99}{Entity-relationship diagram.}

\subsection{\texttt{projects}}

The flagship table. One row per pitch / TED talk / interview. The
\texttt{segment\_type} discriminator (\texttt{'project'},
\texttt{'ted\_talk'}, \texttt{'interview'}) drives the rubric chosen on the
judge UI and the slot a project is shown in on the public gallery. We
denormalise \texttt{vote\_count} onto the row so the public homepage and
leaderboard can be served from a single \texttt{SELECT} ordered by
\texttt{vote\_count}; the count is updated atomically from a Postgres RPC
(see Section~\ref{sub:rpc}).

The \texttt{members} and \texttt{tags} columns are JSONB; they are
displayed but never queried by inner fields. Storing them as JSONB rather
than satellite tables is a deliberate complexity reduction.

\subsection{\texttt{judges} and \texttt{judge\_projects}}

Judges authenticate by email + bcrypt-hashed password. \texttt{is\_admin}
is the role flag; admin judges also see the analytics and integrity
dashboards. The \texttt{judge\_projects} join table is a simple
many-to-many edge between judges and projects; it materialises the
``which projects am I assigned to grade'' query as a single
\texttt{WHERE judge\_id = ?}.

\subsection{\texttt{grades}}

One row per (judge, project) pair, with a unique constraint on that
composite key. The five rubric scores live as separate
\texttt{int} columns rather than JSONB so we can compute aggregates in
SQL (and so the inter-rater reliability calculation can run cleanly in
the controller). The \texttt{status} discriminator
(\texttt{not\_started} / \texttt{in\_progress} / \texttt{submitted}) is
respected in code: once the row's status is \texttt{submitted}, the
controller refuses any further write.

\subsection{\texttt{votes}}

One row per cast vote. The natural identity is the
\texttt{browser\_token} (a UUID generated client-side and persisted in
localStorage); we also record \texttt{ip\_address} for the anomaly
detector. Duplicate prevention is enforced at the application layer
(checking for an existing row with the same browser token); we
deliberately do \emph{not} put a \texttt{UNIQUE} constraint on
\texttt{browser\_token} because a same-browser revote attempt should
return a clean 409 from the controller, not a constraint violation.

\subsection{\texttt{vote\_flags}}

The anomaly detector inserts here when a vote scores at or above the
\texttt{low} threshold. The \texttt{severity} enum is
\texttt{'low' | 'medium' | 'high'}; only \texttt{high} severity
contributes to the People's Choice exclusion. \texttt{reasons} is JSONB
of human-readable strings (e.g.\ \texttt{"burst:5 from same IP within
60s"}) so the integrity dashboard can render them verbatim.

\subsection{\texttt{ai\_insights}}

A per-project cache of the LLM advisory output. Keyed on
\texttt{project\_id} (one-to-one with \texttt{projects}). The
\texttt{model} column captures which model produced this row, so a future
upgrade can selectively re-render older cached entries. The
\texttt{payload} JSONB is the parsed JSON returned by the controller.

\subsection{\texttt{audit\_logs}}

Append-only. Every \texttt{grade.submit}, \texttt{session.set\_public},
\texttt{session.set\_votes\_visible}, and \texttt{session.update} writes a
row here. Indexed on \texttt{action}, \texttt{actor\_id}, and
\texttt{created\_at} for fast forensic replay.

\subsection{\texttt{live\_sessions}}

A single row keyed by \texttt{key = 'main'} that holds the live event's
ephemeral state: \texttt{is\_event\_live} (the master switch),
\texttt{now\_playing} (a JSONB blob with sub-fields
\texttt{voteCountVisible} and \texttt{isPublic} alongside the actual
playing-now segment metadata), and \texttt{up\_next}. We store these as a
single row rather than a key-value store so an admin update is one
\texttt{UPSERT}.

\subsection{\texttt{schedule\_items}}

The fixed event timeline (Welcome, TED Talks, Interviews, Pitches,
Closing). Read at app load to render the timeline strip on the home page;
no writes during the event.

\subsection{Atomic vote increment RPC}
\label{sub:rpc}

The naive way to count votes; read \texttt{vote\_count}, add one,
write back; has a race window. Postgres can do better: a function
defined as

\begin{lstlisting}[language=SQL]
create or replace function increment_vote_count(p_id uuid)
returns void language sql as $$
  update projects set vote_count = vote_count + 1 where id = p_id;
$$;
\end{lstlisting}

is invoked from the controller as
\texttt{await supabase.rpc('increment\_vote\_count', \{ p\_id: projectId \})}.
The atomicity is provided by the database engine; concurrent voters do
not lose updates.

\subsection{Indexes}

Beyond the primary keys and the unique on \texttt{(judge\_id,
project\_id)} in \texttt{grades}, we declare:

\begin{itemize}[leftmargin=*]
  \item \texttt{audit\_logs(action)}, \texttt{audit\_logs(actor\_id)},
        \texttt{audit\_logs(created\_at desc)}; the integrity dashboard
        sorts by recency and filters by action.
  \item \texttt{vote\_flags(project\_id)},
        \texttt{vote\_flags(severity)},
        \texttt{vote\_flags(created\_at desc)}; supports the anomaly
        list and the per-project tainted check.
\end{itemize}

The other read patterns are point-lookups by primary key, which Postgres
indexes automatically.
